% id: 117-03
% book: 베이직쎈 확률과 통계 · 06 이항분포와 정규분포
% section: 개념 35 이항분포와 정규분포의 관계
% passage: 확률변수 $X$가 이항분포 $\mathrm{B}\left(72{,}\mkern3mu\mathopen{}\dfrac{2}{3}\right)$를 따를 때, 오른쪽 표준정규분포표를 이용하여 다음을 구하시오.
% figure: tikz:fig-117-03
% answer: 
% answer_source: 
% variant_level: 0 (원본 전사)
% 이 파일은 build.mjs 가 items.json 에서 만든다. 고칠 때는 items.json 을 고친다.
\smash{\raisebox{1.5mm}{\dmkichul{베이직쎈 확률과 통계 117쪽 03번}}}
$\mathrm{P}(48\le X\le 56)$
\dmhint{$\mathrm{E}(X)=72\cdot\blank=\blank$, $\mathrm{V}(X)=72\cdot\dfrac{2}{3}\cdot\blank=\blank$\\ 이때 72는 충분히 큰 수이므로 $X$는 근사적으로 정규분포 $\mathrm{N}(\blank{,}\mkern3mu\mathopen{}{\blank}^2)$을 따른다.\\ 따라서 $Z=\dfrac{X-\blank}{\blank}$로 놓으면 확률변수 $Z$는 표준정규분포 $\mathrm{N}(0{,}\mkern3mu\mathopen{}1)$을 따르므로\\ $\mathrm{P}(48\le X\le 56)=\mathrm{P}(\blank\le Z\le\blank)=\blank[2.4em]$}
